\par\Needspace{8\baselineskip}\bigskip
\SourceLeaf{487}%
\pdfbookmark[0]{Plan for the degree survey. May 30, 1819.}{entwurf-gradmessung-mai1819}
\begin{center}\large [PLAN FOR THE DEGREE SURVEY.]\end{center}
\par
— — — In a letter to Your — —, I need say nothing about degree surveys in general or the extent to which they constitute one of the most interesting subjects of human knowledge. The great benefits and the new light on still obscure points that may be expected from carrying out more such undertakings depend in part on their being conducted on as large a scale as possible. Isolated degree surveys in Europe covering only a small extent could now, after the great works in France and England, have only a very subordinate value. One or a couple of further surveys of this kind in Europe, over a considerable extent, would certainly be of immense importance for our knowledge of the shape of the Earth.

A glance at the map of Europe and at the state of civilization of its various nations shows that, apart from the great line from the Balearic Islands to the Orkney Islands, similar operations will be feasible in only two more directions, 1) in the Russian Empire and 2) through the Jutland peninsula and the whole of Germany as far as the Mediterranean. According to reports that have reached me, there appear to be some prospects of carrying out a survey in the former empire. For the other, however, there is more than a prospect. The first major step has already actually been taken. The degree survey ordered by the King of Denmark, from the northern tip of Jutland to Lauenburg, has already been under way for two years. The skill and judgment of the Danish astronomer, the excellence of his instruments and the generosity with which the King of Denmark approves everything found necessary or desirable for the perfection of the survey are a guarantee that it will be carried out exactly as the present state of the sciences and arts makes possible and necessary.

The Danish degree survey is to include, in addition to the meridian arc mentioned above, the measurement of an arc of a parallel of latitude from the west coast of Jutland to Copenhagen. Naturally, only the first part of the plan can be discussed here. That arc by itself already spans \(4\frac12\) degrees, and the survey is therefore of respectable extent even when considered on its own. But it appears incomparably more important when regarded as the beginning of that great survey which could be extended as far as the island of Elba, and thus to \(16^0\). It is probably more than a fanciful hope that sooner or later this undertaking will be carried out to such an extent, especially since some preliminary work has in fact already been done in a few of the intervening lands, particularly in Thuringia and Bavaria. But the continuation of the Danish survey through the Kingdom of Hanover is the first and most essential condition for the eventual realization of that plan. This continuation alone would extend the arc by 2 degrees. If the Gotha observatory were then also linked to the Göttingen observatory by a network of triangles, which would be very desirable in several other respects as well and would be easy to carry out, an arc of 7 degrees would already be realized.

I need only mention briefly that the measurement of a meridian arc from Hamburg to Göttingen would be of great importance in respects other than the purely scientific. If sooner or later it were decided to survey the whole Kingdom of Hanover in accordance with present requirements, the network of triangles established for this purpose would provide the surest foundation for connecting further triangulations to it eastward and westward. And if there were a near prospect of such a general survey, the degree survey could also serve the additional purpose of helping to prepare suitable people for that work.
